\section{Experimental details}
\label{sec:appendix_exp_detail}
\subsection{Pretraining on OpenWebText}
\paragraph{Dataset preparation.}
As mentioned in Section~\ref{sec:experiment_pretrain}, to obtain a variable-length dataset, we split OpenWebText articles paragraph-wise using the GPT-2 tokenizer~\cite{radford2019rewon}. This can be implemented by locating the token index corresponding to the delimiter $\verb|\n |\verb|\n|$. Sequences longer than $1024$ tokens are then chunked recursively by splitting by the delimiter closest to the middle sequence, yielding a variable-length dataset with maximum sequence length $1024$.

\paragraph{FlexAttention.}
To handle variable-length sequences during training, we pad each batch to the maximum sequence length. In MDM, by design, pad tokens also enter the model input. In contrast, FlexMDM is designed to receive only clean or mask tokens as inputs. Ideally, QKV attention should not attend to pad tokens; however, current FlashAttention~\citep{dao2022flashattention} does not support this for \emph{non-causal} attention (our setup of interest). We therefore adapt FlexAttention~\citep{dong2024flex}. A side benefit is improved training speed, since FlexMDM performs attention over fewer tokens than MDM’s full-sequence attention. Note that in the LLaDA experiment, we did not implement this optimization; pad tokens can therefore attend to other tokens, though we expect the impact to be negligible.

\paragraph{Training configuration.}
As in Section~\ref{sec:experiment_pretrain}, we model FlexMDM with a DiT~\cite{peebles2023scalable} backbone and embed the insertion schedule $\alpha_t$. For MDM, we use the same DiT configuration with time step embedding but without the softplus scalar head. Transformer configuration is: \texttt{hidden size:768}, \texttt{heads:12}, \texttt{conditional dimension:128}, \texttt{dropout:0.05}, \texttt{layers:12}.
We train both models on $16$ H100 GPUs with a global batch size of $1024$ and max training iteration $1M$. We use the AdamW~\citep{loshchilov2017decoupled} optimizer with weight decay $0.03$, learning rate $3\mathrm{e}{-4}$, $2000$ warmup steps, and an EMA factor of $0.9999$. Additionally, we use low-discrepancy time-step sampling to reduce variance: one $t$ is drawn uniformly from each interval ${[i/T,(i+1)/T]}_{i=0}^{T-1}$, as in prior MDM training~\citep{shi2024simplified}.

\paragraph{Metric.}
For evaluation, we take sequences generated by both models and retain the clean tokens by removing padding (e.g., the leading pad token). We adopt LLaMA-2-7B~\citep{touvron2023llama2} as the reference model to compute likelihoods. We notice that the pretrained MDM generates short sentences with unreasonably large (worse) perplexities.  Therefore, we filter overly short sequences with $\le$ 10 tokens when calculating average perplexity.

\subsection{Pretraining on the Maze planning task}
\label{app:maze}
\paragraph{Task construction.}
We generate mazes with a fully random, recursive division procedure (the code is provided in Listing~\ref{lst:example}), on a $41\times 41$ grid, with some invalid cells. As described in Section~\ref{sec:experiment_pretrain}, we consider a subgoal-conditioned planning task: the model is given a sequence of subgoals $(g_1,\dots,g_K)$ and must produce a valid path that connects them in order. To construct training examples for a given $K$, we sample $15000$ start–goal pairs ($g_1,g_K$), compute the shortest path for each pair via breadth-first search, and then add controlled detours to obtain up to $10$ distinct valid paths per pair. Subgoals are formed by selecting $K-2$ intermediate cells uniformly at random along a chosen path (start and goal are already fixed). We use $95\%$ of the pairs for training and hold out $5\%$ for validation to evaluate generalization to unseen pairs and subgoal sets.

\paragraph{Training data construction.}
For MDM, the training sequence is $((g_1,\dots,g_K)\;\texttt{[SEP]}\;\texttt{Path})$,
where \texttt{[SEP]} is a special separator and \texttt{Path} denotes the tokenized path. During training, the prompt $(g_1,\dots,g_K)$ bypasses the interpolant so that, at inference time, the model can condition on $(g_1,\dots,g_K)\;\texttt{[SEP]}$ and generate the path. For FlexMDM, we use an interpolant in which the subgoal tokens are exempt from the process in~\eqref{eqn::FlexMDM_interpolant}; that is, tokens corresponding to each $g_i$ are kept clean at all times. This also enables generation to start from $(g_1,\dots,g_K)$. Although this conditional generation template changes the base distributions $p_0$ for both MDM and FlexMDM, we note that the theoretical guarantee from Section~\ref{sec:FlexMDM_informal} and Section~\ref{sec:FlexMDM_inference} still holds--once the training is perfect (under the access to the ground truth unmasking posterior and insertion expectation), the inference algorithms recover the ground truth distribution $p_1$.

\paragraph{Training configuration.}
We use the same architectural design as in the OpenWebText pretraining, but with a smaller model:
\texttt{hidden size:256}, \texttt{heads:8}, \texttt{conditional dimension:128}, \texttt{dropout:0.1}, \texttt{layers:8}.
Both models have approximately $8.5$M parameters. We train them on $4$ A100 GPUs with a global batch size of $256$ for up to $100$ epochs. We use AdamW~\citep{loshchilov2017decoupled} with weight decay $0.01$, learning rate $1\times 10^{-4}$, $20$ warmup epochs, and an EMA factor of $0.9999$.

\paragraph{Metric.}
Given the final conditionally generated sequence, we report the success rate under two criteria: (1) all visited cells are valid (no collisions with invalid cells), and (2) the path connects the subgoals consecutively in order. We perform inference both models with the number of sampling steps $256$.


\lstset{
  language=Python,
  basicstyle=\ttfamily\small,
  keywordstyle=\color{xblue},
  stringstyle=\color{xxgreen},
  commentstyle=\color{gray},
  showstringspaces=false,
  breaklines=true
}
\begin{center}
\begin{lstlisting}[caption={Code for the maze Construction}, label={lst:example}]
# ----------------------------------------------------------------------
#  RECURSIVE DIVISION (perfect maze) -----------------------------------
# ----------------------------------------------------------------------
def _divide(g, top, left, h, w):
    if h <= 2 or w <= 2:
        return
    horizontal = w < h  # split the longer dimension
    if horizontal:
        y = random.randrange(top + 1, top + h - 1, 2)      
        gap = random.randrange(left, left + w, 2)          
        g[y, left:left + w] = 1
        g[y, gap] = 0                                      
        _divide(g, top, left, y - top, w)     
        _divide(g, y + 1, left, top + h - y - 1, w)    
    else:
        x = random.randrange(left + 1, left + w - 1, 2)   
        gap = random.randrange(top, top + h, 2)
        g[top:top + h, x] = 1
        g[gap, x] = 0
        _divide(g, top, left, h, x - left) 
        _divide(g, top, x + 1, h, left + w - x - 1)


# ----------------------------------------------------------------------
#  WRAPPER with extra doorways ----------------------
# ----------------------------------------------------------------------
def division_maze(m, n, seed=2025, extra_door_frac=0.5):
    """
    m, n             # size in *cells*      (not bitmap coords)
    seed             # int or None
    extra_door_frac  # 0, perfect maze; >0 flicks more doors (loops)
    """
    random.seed(seed)
    H, W = 2 * m + 1, 2 * n + 1
    g = np.zeros((H, W), dtype=np.uint8)
    g[0, :], g[H - 1, :], g[:, 0], g[:, W - 1] = 1, 1, 1, 1

    _divide(g, 1, 1, H - 2, W - 2)

    # ---------- optional imperfection: punch extra doorways ------------
    if extra_door_frac > 0:
        candidates = []
        for r in range(1, H - 1):
            for c in range(1, W - 1):
                if g[r, c] == 1:
                    # Check if wall separates two passages orthogonally
                    if g[r - 1, c] == 0 and g[r + 1, c] == 0:
                        candidates.append((r, c))
                    elif g[r, c - 1] == 0 and g[r, c + 1] == 0:
                        candidates.append((r, c))
        k = int(len(candidates) * extra_door_frac)
        for (r, c) in random.sample(candidates, k=k):
            g[r, c] = 0
    return g
\end{lstlisting}
\end{center}



\subsection{Weight Initialization training from LLaDA}
In this section, we describe the procedure for adapting the pretrained LLaDA-8B base model into the FlexMDM.

\paragraph{Training configuration.}
LLaDA uses a bidirectional Transformer without an additional time embedding, leveraging the fact that MDM does not require an explicit time signal \citep{zheng2024masked}. For FlexMDM, to model the insertion expectation, we inject a time-embedding pathway via AdaLN \citep{peebles2023scalable}. For parameter efficiency, we tie (share) the four AdaLN parameter sets across the intermediate Transformer layers. We also attach a softplus scalar head to parameterize the insertion expectation.

Next, we attach LoRA adapters \citep{hu2022lora} to every attention projection (q\_proj, k\_proj, v\_proj) and to the unmasking-posterior head. We include LoRA on the unmasking posterior head because the unmasking posteriors differ between MDM and FlexMDM: in FlexMDM, unmasking must account for token shifts induced by insertions. This fine-tuning of the last head is crucial for enabling variable-length generation.

The LoRA configuration that we use is $r=128$, $\alpha=128$, and dropout $0.1$. Training runs for $200,000$ gradient steps with a batch size of $64$ on $16$ H100s, which took approximately 3 days. We optimize with AdamW \citep{loshchilov2017decoupled} at learning rate $1\times10^{-4}$ and weight decay $0.1$, using a cosine warm-restarts scheduler.

\paragraph{Evaluation on GSM8K.}
We instruction-fine-tune (IFT) the FlexMDM base model on the GSM8K training split. To preserve the instruction–answer format, we modify the interpolant in Eq.~\ref{eqn::FlexMDM_interpolant} so that tokens corresponding to the instruction are excluded from the interpolant—these tokens remain fixed for all time steps. We apply the same strategy to obtain the baseline (that is, IFT from LLaDA-base), modifying the MDM interpolant so that instruction tokens remain fixed. This choice is equivalent to the IFT recipe used in \cite{nie2025large,dream2025}. Both models are IFT-ed for $1000$ epochs. (Other IFT hyperparameters match those used in our base setup unless otherwise noted.)

For FlexMDM inference, we start from the instruction prompt at $t=0$ and run adaptive inference to $t=1$. Concretely, we use Top-K probability with a sliding window (Appendix~\ref{sec:appendix_theory_inference}) with $\gamma_1=5.0$ and $\gamma_2=64$. For LLaDA, we report the best result under the semi-autoregressive, confidence-based sampling of \citep{nie2025large}. For both models, we set the token sampling temperature to $0.0$, which we confirm to be important for strong Pass@1. Overall, adaptive inference substantially improves performance over vanilla inference.

\paragraph{Evaluation on HumanEval-infill.}
Code infilling conditions on a prefix and suffix, and asks the model to complete the middle part of the code. For FlexMDM, we format training examples as
$(\texttt{Prefix};\texttt{[SEP]};\texttt{[Middle]};\texttt{[SEP]};\texttt{Suffix})$,
where \texttt{[SEP]} is a separator token and \texttt{Instruction} describes the infilling task for a model. We modify Eq.~\eqref{eqn::FlexMDM_interpolant} so that tokens for \texttt{Prefix}, \texttt{Suffix}, and \texttt{[SEP]} are fixed in the interpolant (i.e., excluded from the interpolant). Thus, at $t=0$ the state is
$(\texttt{Prefix};\texttt{[SEP]};\texttt{[SEP]};\texttt{Suffix})$.

For MDM, we use the format
$(\texttt{Instruction};\texttt{[PRE]};\texttt{Prefix};\texttt{[SUF ]};\texttt{Suffix};\texttt{[SEP]};\texttt{Middle})$,
with \texttt{Instruction} prompting infill after prefix and suffix
, along with \texttt{[PRE]} and \texttt{[SUF]} separate the prefix and suffix, respectively. Here too, the tokens without \texttt{Middle} are held fixed by the modified interpolant. This difference in formatting reflects the fixed-length nature of MDMf MDM (no token insertion). This formatting has been used in the code infilling tasks for autoregressive models in \cite{bavarian2022efficient}. Naively masking the \texttt{Middle} span yields a base state at $t=0$ of $(\texttt{Prefix};\texttt{Masked};\texttt{Suffix})$, where \texttt{Masked} is a fully masked sequence of length $|\texttt{Middle}|$. This leaks length information—materially simplifying the task— and renders comparisons to FlexMDM unfair, since FlexMDM does not observe the target span length. For fair evaluation, we therefore avoid length-revealing masks and require methods to infer the span length during inference.

We IFT both models on the educational-instruct split of \texttt{opc-sft-stage2}; the architecture and optimization configurations match those used for GSM8K IFT. At evaluation, we initialize from the base distributions: for FlexMDM,
$(\texttt{Instruction};\texttt{Prefix};\texttt{[SEP]};\texttt{[SEP]};\texttt{Suffix})$;
for MDM,
$(\texttt{Instruction};\texttt{Prefix};\texttt{Suffix};\texttt{[SEP]})$.
We use the same Top-K adaptive inference for both and temperature $0.0$. Final outputs are scored with the HumanEval-infill verifier toolkit to compute success rates.